\section{From formulas to bounded-degree labelled DAGs}\label{sec:formula-dag}

We recall a standard syntax-directed conversion.  Its precise form is useful because the
support and counting arguments depend on local degrees.

\begin{construction}[Formula DAG]\label{cons:dag}
Associate to each formula \(\Phi\) a labelled ranked DAG \(G_\Phi\) as follows.
\begin{enumerate}[label=(\roman*)]
\item A leaf \(\ell\) becomes one edge from a fresh source to a fresh sink, labelled
\(\ell\).
\item If \(\Phi=\Psi\Theta\), take the disjoint union of \(G_\Psi\) and \(G_\Theta\) and
add one unit-labelled edge from the sink of \(G_\Psi\) to the source of \(G_\Theta\).
\item If \(\Phi=\Psi+\Theta\), take the disjoint union, add a fresh source with one
unit-labelled edge to each child source, and add one unit-labelled edge from each child
sink to a fresh sink.
\end{enumerate}
Ranks are chosen compatibly so every edge increases rank.
\end{construction}

\begin{lemma}[Value and local degree]\label{lem:dag-properties}
The path polynomial of \(G_\Phi\) equals \(\Phi\).  Every edge label is a variable or a
constant.  The source has indegree zero and outdegree at most two; the sink has outdegree
zero and indegree at most two; every internal vertex has total degree at most three.
\end{lemma}

\begin{proof}
The value statement follows by structural induction.  A leaf has its label as its unique
path value.  The multiplication construction concatenates a path in the left graph with
the unit connector and a path in the right graph, multiplying their values.  The addition
construction offers a disjoint choice between the two child graphs, adding their path
sums.  The degree assertions follow simultaneously from the three constructors.  In the
multiplication constructor, the old left sink gains one outgoing connector and the old
right source gains one incoming connector; each had the opposite boundary degree zero.
In the addition constructor, the new boundary vertices have degree two, while each child
boundary gains one incident connector.  Thus no internal total degree exceeds three.
\end{proof}

Let \(l,a,u\) be the leaf, addition-gate, and multiplication-gate counts of \(\Phi\).

\begin{lemma}[Graph counts]\label{lem:graph-counts}
For \(G_\Phi=(V,E)\),
\[
  |E|=l+4a+u,
  \qquad
  |V|=2l+2a.
\]
\end{lemma}

\begin{proof}
Induct on the formula.  A leaf contributes two vertices and one edge.  A multiplication
adds the counts of the children and one connector edge.  An addition adds the child counts,
two new vertices, and four connector edges.
\end{proof}

\section{The Boolean path selector}\label{sec:selector}

Fix a ranked DAG \(G=(V,E)\) satisfying the degree conditions of
\cref{lem:dag-properties}.  Introduce a Boolean variable \(z_e\) for each edge.  For a
vertex \(v\), set
\[
 I_v(z)=\sum_{e:\,\mathrm{tgt}(e)=v}z_e,
 \qquad
 O_v(z)=\sum_{e:\,\mathrm{src}(e)=v}z_e.
\]
Define the flow factor
\[
 F_v(z)=
 \begin{cases}
   O_s(z),&v=s,\\
   I_t(z),&v=t,\\
   1+I_v(z)+O_v(z),&v\notin\{s,t\}.
 \end{cases}
\]
For every ordered pair of distinct edges with a common head, and every ordered pair of
distinct edges with a common tail, include a pair-exclusion factor \(1+z_ez_{e'}\).  Let
\(\mathcal P\) be the resulting ordered-pair index set, and put
\[
 \Sel_G(z)=
 \left(\prod_{v\in V}F_v(z)\right)
 \left(\prod_{(e,e')\in\mathcal P}(1+z_ez_{e'})\right).
\]
The use of ordered pairs duplicates each unordered exclusion.  This does not change its
Boolean value, which is zero or one.

\begin{proposition}[Selector correctness]\label{prop:selector}
For every \(z\in\B^E\),
\[
 \Sel_G(z)=
 \begin{cases}
  1,&\{e:z_e=1\}\text{ is the edge set of a directed }s\text{--}t\text{ path},\\
  0,&\text{otherwise}.
 \end{cases}
\]
\end{proposition}

\begin{proof}
Every factor evaluates to zero or one.  The pair exclusions force at most one selected
incoming edge and at most one selected outgoing edge at each vertex.  Subject to those
exclusions, the source flow factor is one precisely when the source has one selected
outgoing edge, and the sink factor is one precisely when the sink has one selected incoming
edge.  At an internal vertex, \(I_v,O_v\in\{0,1\}\), and
\(1+I_v+O_v=1\) precisely when \(I_v=O_v\).  Hence every selected internal vertex has one
selected incoming edge if and only if it has one selected outgoing edge.

Starting at the unique selected edge out of \(s\), conservation and finiteness extend it
forward until it reaches \(t\).  It cannot stop at an internal vertex.  It cannot revisit a
vertex because the graph is ranked and acyclic.  Thus it gives an \(s\)-\(t\) path.  Any
additional selected component would be finite, avoid \(s,t\), and have equal selected
indegree and outdegree at every one of its vertices.  Following its unique outgoing edges
would create a directed cycle, contradicting acyclicity.  Therefore no additional component
exists.  The converse is immediate: an \(s\)-\(t\) path satisfies all flow and exclusion
factors.
\end{proof}

Let \(\ell_e\) be the affine label of edge \(e\), and define its conditional label factor
\[
  C_e(z_e)=1+z_e(\ell_e+1).
\]
In characteristic two,
\[
  C_e(0)=1,
  \qquad
  C_e(1)=\ell_e.
\]

\begin{proposition}[ABP hypercube identity]\label{prop:abp-cube}
For a labelled DAG over a field of characteristic two,
\[
  \val(G)=
  \sum_{z\in\B^E}\Sel_G(z)\prod_{e\in E}C_e(z_e).
\]
\end{proposition}

\begin{proof}
By \cref{prop:selector}, only selector assignments that are edge sets of source--sink paths
survive.  For such an assignment, the conditional label product equals the product of the
labels on the selected path; all unselected edges contribute one.  Summing the surviving
terms is exactly the path polynomial.
\end{proof}

\section{The support-two compiler}\label{sec:compiler}

The expression in \cref{prop:abp-cube} is not yet a product of support-two affine forms.
It contains quadratic factors \(1+xy\) and, at total-degree-three internal vertices,
ternary affine flow factors.  We now compile both kinds with fresh one-bit gadgets.

\subsection{Occurrence sets and fresh auxiliaries}

Let
\[
  D=\#\{v\in V\setminus\{s,t\}:\indeg(v)+\outdeg(v)=3\}
\]
and let \(P=|\mathcal P|\) be the ordered pair count.  There are exactly
\(|E|+P\) quadratic occurrences: one conditional edge factor per edge and one exclusion
per ordered pair.  Give every such occurrence its own fresh bit.  Give every degree-three
internal flow factor its own fresh bit.  The full auxiliary set is the disjoint union of
\[
  \underbrace{E}_{\text{original selectors}}
  \ \sqcup\ 
  \underbrace{T}_{\text{ternary gadgets}}
  \ \sqcup\ 
  \underbrace{K}_{\text{quadratic gadgets}},
\]
where \(|T|=D\) and \(|K|=|E|+P\).

Freshness is structural.  It prevents the generally false replacement
\[
  \left(\sum_h A_h\right)\left(\sum_k B_k\right)
  \stackrel{\text{false}}{=}
  \sum_h A_hB_h.
\]
With disjoint bits, finite distributivity gives a sum over the Cartesian product of the
local cubes.

\subsection{Compiling quadratic occurrences}

Every pair exclusion is \(1+z_ez_{e'}\).  Every conditional edge factor can be written
\[
  1+z_e(\ell_e+1).
\]
In a formula-generated DAG, \(\ell_e\) is a variable or constant.  Therefore each of the
two arguments of the quadratic gadget in \cref{lem:quadratic} has support at most one.
Replacing each quadratic occurrence by its own copy of the gadget yields three affine
factors of support at most two.

\subsection{Compiling ternary flow factors}

At a degree-three internal vertex, the flow factor is
\[
  1+z_a+z_b+z_c
\]
for its three incident edges.  Replace it by the four factors from
\cref{lem:ternary}.  After summing the fresh ternary bit, the resulting local value is
\[
  1+z_a+z_b+z_c+\kappa^{-1}z_az_bz_c.
\]
The last summand is the error term.

Among three incident edges and two sides of a vertex, two edges lie on the same side.
Choose \(a,b\) so that they share a head or share a tail.  The original selector contains
the corresponding exclusion \(1+z_az_b\).

\begin{lemma}[Pointwise ternary cancellation]\label{lem:local-cancel}
Fix an original selector assignment \(z\in\B^E\).  Let \(e_v(z)=\kappa^{-1}z_az_bz_c\) be
the error at a ternary vertex \(v\), and let \(W(z)\) be the product of all recovered
quadratic factors.  Then
\[
  e_v(z)W(z)=0.
\]
\end{lemma}

\begin{proof}
The product \(W(z)\) contains \(1+z_az_b\).  Hence the claim follows from
\cref{lem:killer} after multiplying by the remaining factors.
\end{proof}

The following elementary product lemma lets us remove all errors simultaneously.

\begin{lemma}[Product error elimination]\label{lem:prod-errors}
Let \(R\) be a commutative ring, let \(p_j,e_j,W\in R\), and suppose
\(e_jW=0\) for every \(j\) in a finite index set.  Then
\[
  \left(\prod_j(p_j+e_j)\right)W
  =\left(\prod_jp_j\right)W.
\]
\end{lemma}

\begin{proof}
Induct on the index set.  In the induction step,
\[
 (p+e)XW=pXW+eXW=pXW,
\]
because commutativity gives \(eXW=X(eW)=0\).
\end{proof}

\begin{theorem}[Graph compiler]\label{thm:graph-compiler}
Let \(G\) satisfy the degree and label conditions of \cref{lem:dag-properties}.  Over a
characteristic-two field containing \(\tau\notin\{0,1\}\), the path polynomial of \(G\)
has a support-two affine-hypercube representation with
\[
 q=2|E|+D+P,
 \qquad
 M=|V|+4D+3(|E|+P).
\]
\end{theorem}

\begin{proof}
Fix the original selector assignment \(z\).  Sum first over all fresh quadratic and
ternary gadget bits.  Because every occurrence has a distinct fresh auxiliary, repeated
finite distributivity factors these local sums.  By \cref{lem:quadratic}, the quadratic
sums recover every conditional edge factor and pair exclusion.  By
\cref{lem:ternary}, each ternary sum recovers the corresponding flow factor plus its cubic
error.  Apply \cref{lem:local-cancel,lem:prod-errors} while \(z\) is fixed.  The resulting
expression is precisely
\[
  \Sel_G(z)\prod_{e\in E}C_e(z_e).
\]
Now sum over \(z\) and invoke \cref{prop:abp-cube}.

For support, the source and sink flow factors use at most two selectors.  A non-ternary
internal flow factor uses at most two selectors, and a ternary vertex contributes a
constant placeholder together with four support-two gadget factors.  Every quadratic
gadget factor has support at most two.  The auxiliary count is
\[
 |E|+D+(|E|+P)=2|E|+D+P.
\]
The implemented factor list contains one flow placeholder for every vertex, four factors
per ternary occurrence, and three factors per quadratic occurrence, giving
\[
 |V|+4D+3(|E|+P).
\]
\end{proof}

\begin{remark}
The constant placeholder at a ternary vertex is retained in the formal representation.
Dropping such unit factors could improve numerical constants, but is irrelevant to the
class theorem and is not used in the certified accounting.
\end{remark}
